\documentclass[11pt,a4paper]{article}
\usepackage[utf8]{inputenc}
\usepackage[T1]{fontenc}
\usepackage{amsmath,amssymb,amsthm,mathtools}
\usepackage{geometry}
\geometry{margin=1in}
\usepackage{graphicx}
\usepackage{float}
\usepackage[hidelinks]{hyperref}
\usepackage{cleveref}
\usepackage{booktabs}
\usepackage{array}
\usepackage{enumitem}
\usepackage{caption}
\usepackage{subcaption}
\usepackage{algorithm}
\usepackage{algpseudocode}
\usepackage[numbers]{natbib}   % [numbers] added: plain style -> numeric citations
\bibliographystyle{plainnat}
\newtheorem{theorem}{Theorem}[section]
\newtheorem{lemma}[theorem]{Lemma}
\newtheorem{proposition}[theorem]{Proposition}
\newtheorem{corollary}[theorem]{Corollary}
\theoremstyle{definition}
\newtheorem{definition}[theorem]{Definition}
\newtheorem{remark}[theorem]{Remark}
\newcommand{\R}{\mathbb{R}}
\newcommand{\Z}{\mathbb{Z}}
\newcommand{\N}{\mathbb{N}}
\newcommand{\E}{\mathbb{E}}
\newcommand{\Var}{\mathrm{Var}}
\newcommand{\Cov}{\mathrm{Cov}}
\newcommand{\tr}{\mathrm{tr}}
\newcommand{\dd}{\,\mathrm{d}}
\newcommand{\avg}[1]{\langle #1 \rangle}
% --- added: macros used in the body but absent from the provided preamble ---
\numberwithin{equation}{section}   % added: per-section equation numbers (matches theorem[section])
\newcommand{\bI}{\mathbf{I}}
\newcommand{\bphi}{\boldsymbol{\varphi}}
\newcommand{\Lgen}{\mathcal{L}}
\newcommand{\Ladj}{\mathcal{L}^{*}}

\title{Numerical study of an energy cascade model derived from a dispersive equation}
\author{Jayleen Jiang \and Yao Li\thanks{Department of Mathematics and Statistics, University of Massachusetts Amherst, Amherst, MA 01003, USA. \texttt{liyao@umass.edu}}}
\date{}

\begin{document}
\maketitle

\begin{abstract}
We study a stochastic energy-cascade chain obtained from the resonant toy model
for the cubic nonlinear Schr\"odinger equation. The deterministic chain
transports action between generations of resonant Fourier modes; coupling the
two endpoints to heat baths produces a nonequilibrium steady state (NESS) whose
long-chain behavior is analytically open. We use a Gibbs-preserving
fluctuation--dissipation bath, validate the resulting equilibrium law against
an independent Gibbs sampler, and compute the NESS by high-throughput Monte
Carlo and mesh-free Fokker--Planck solvers. For closed chains with
$T_1=10$, $T_n=2$, and $\gamma=0.1$, the mean conserved action current obeys
$\E[J(n)]=28.75\,n^{-1.850}$ over $n=10,20,30,40$, with bootstrap exponent
confidence interval $[-1.870,-1.830]$. Thus the finite-chain action
conductivity scales as $\kappa_n\sim n^{-0.850}$, indicating an insulating
action-transport regime rather than Fourier transport. Long-chain pair
marginals are well described by rescaled equilibrium marginals at the local
kinetic temperature, but strict local Gibbs structure fails even at
equilibrium because of the nonlocal mass term in the Hamiltonian. For the
three-mode chain we validate a neural-network Fokker--Planck solver against the
explicit Gibbs density, then use it to visualize symmetry breaking,
phase-locking stabilization, and qualitative slow relaxation modes. The
transport and local-equilibrium conclusions are supported by reproducible
simulation code and finite-time diagnostics; bath-energy entropy production and
asymptotic current large deviations are left for future work.
\end{abstract}


\section{Introduction}

The transfer of energy between the modes of a nonlinear Hamiltonian system is a
basic question across the sciences, from fluids and oceanography to climate.
When external forcing and dissipation act on a few modes while the bulk---the
\emph{inertial range}---evolves intrinsically, the natural framework is
out-of-equilibrium statistical physics; for wave systems it is \emph{wave} (or
\emph{weak}) \emph{turbulence theory} \cite{Nazarenko}, in which energy is
transported between Fourier modes, that is, across spatial scales. The classical
description proceeds through the wave kinetic equation, whose power-law
stationary solutions are the Kolmogorov--Zakharov cascade spectra
$n(k)=|k|^{-d}$ and $n(k)=|k|^{-d+2/3}$ (the direct and inverse cascades). These
spectra are themselves nonequilibrium steady states (NESS), in which energy
flows at constant flux through the inertial range; their rigorous justification
remains largely open \cite{HLNS}.

Following \cite{HLNS}, we adopt a complementary viewpoint: rather than the
kinetic equation, we study the invariant measure---the NESS itself---of a
forced and dissipated reduction of the nonlinear Schr\"odinger (NLS) equation.
The present paper is the numerical companion to that program. We compute and
analyze the NESS of the energy cascade chain: for the three-mode system treated
analytically in \cite{HLNS}, and, for the first time, for long chains, where the
analysis is open.

\subsection{Derivation of the model}\label{sec:derivation}

Consider the cubic NLS on a torus,
\begin{equation}\label{eq:nls}
  -i\partial_t v + \Delta v = \pm |v|^2 v .
\end{equation}
In Fourier variables $a_K$ the cubic nonlinearity couples wavevectors with
$K_1-K_2+K_3=K$. Retaining only the \emph{resonant} interactions, those with
$|K_1|^2-|K_2|^2+|K_3|^2=|K|^2$, gives the resonant NLS system, which
approximates \eqref{eq:nls} for sufficiently small data. Colliander, Keel,
Staffilani, Takaoka and Tao \cite{CKSTT} restrict this system to a finite set of
frequencies
\[
  \Lambda=\Lambda_1\cup\Lambda_2\cup\cdots\cup\Lambda_n\subset\mathbb{Z}^2,
\]
on which the amplitude is constant within each generation $\Lambda_j$. The
generations are produced by an explicit combinatorial ``nuclear family''
construction that embeds successive $\Lambda_j$ in $\mathbb{Z}^2$ through exact
rational (Pythagorean) rotations, with $\Lambda_1$ supported at low and
$\Lambda_n$ at high frequency. Writing $c_j(t)$ for the common amplitude on
$\Lambda_j$, the resonant system collapses to the \emph{toy model} of
\cite{CKSTT},
\begin{equation}\label{eq:toy}
  -i\dot c_j = 2\Big(\sum_k |c_k|^2\Big)c_j - |c_j|^2 c_j
  + 2\big(c_{j-1}^2 + c_{j+1}^2\big)\bar c_j,
  \qquad j=1,\dots,n,\quad c_0=c_{n+1}=0 .
\end{equation}
In \cite{CKSTT} this model is shown to admit an energy-cascade solution whose
mass, concentrated on $\Lambda_1$ at $t=0$, migrates to $\Lambda_n$ at a later
time---a deterministic route to the growth of Sobolev norms of \eqref{eq:nls}.

System \eqref{eq:toy} is Hamiltonian. In action--angle coordinates
$c_j=\sqrt{I_j}\,e^{i\varphi_j}$, with the \emph{bond angles}
$\theta_k:=2(\varphi_k-\varphi_{k-1})$, it is generated by the Hamiltonian
\begin{equation}\label{eq:hamiltonian}
  H(\bI,\bphi) \;=\; \Big(\sum_{k=1}^{n} I_k\Big)^{2}
  \;-\;\tfrac12\sum_{k=1}^{n} I_k^{2}
  \;+\; 2\sum_{k=2}^{n} I_{k-1} I_k \cos\theta_k ,
\end{equation}
with the Hamiltonian equations of motion
\begin{equation}\label{eq:ham-flow}
  -\partial_{\varphi_j}H = 4 I_j\big(I_{j-1}\sin\theta_j - I_{j+1}\sin\theta_{j+1}\big),
  \qquad
  \partial_{I_j}H = 2M - I_j + 2 I_{j-1}\cos\theta_j + 2 I_{j+1}\cos\theta_{j+1},
\end{equation}
where $M=\sum_k I_k$ and $I_0=I_{n+1}=0$.

\begin{remark}[Normalization convention]\label{rem:convention}
The Hamiltonian \eqref{eq:hamiltonian} is twice the action--angle Hamiltonian of
\cite{HLNS}; equivalently the drift \eqref{eq:ham-flow} is twice theirs, and the
equilibrium measure below carries the factor $1/2T$ rather than $1/T$. This
amounts to a rescaling of time and is fixed throughout so that every formula
matches the simulation code that produces our numerical data. For $n=3$ with the
middle mode as reference one writes $\theta_1=2(\varphi_1-\varphi_2)$,
$\theta_3=2(\varphi_3-\varphi_2)$, recovering the reduced coordinates of
\cite{HLNS}.
\end{remark}

To realize an energy cascade as a statistical steady state one injects energy at
the low-frequency end and extracts it at the high-frequency end: the two end
modes $j=1$ and $j=n$ are coupled to stochastic heat baths at temperatures $T_1$
and $T_n$. This produces the It\^o system
\begin{align}\label{eq:sde-bulk}
  \dd I_j   &= -\,\partial_{\varphi_j} H \dd t
             \;+\; b^{I}_j \dd t \;+\; \sigma^{I}_j \dd W^{I}_j, \\
  \dd \varphi_j &= \phantom{-}\partial_{I_j} H \dd t
             \;+\; b^{\varphi}_j \dd t \;+\; \sigma^{\varphi}_j \dd W^{\varphi}_j,
             \qquad j = 1,\dots,n, \nonumber
\end{align}
in which the bath drift and diffusion $(b,\sigma)$ are supported only on
$j\in\{1,n\}$. We take the fluctuation--dissipation (Langevin) form: at the hot
end $j=1$,
\begin{equation}\label{eq:bath}
  b^{I}_1 = 2\gamma\big(2T_1 - I_1\,\partial_{I_1}H\big),\quad
  \sigma^{I}_1 = 2\sqrt{2\gamma T_1 I_1},\quad
  b^{\varphi}_1 = 2\gamma I_2 \sin\theta_2,\quad
  \sigma^{\varphi}_1 = \sqrt{2\gamma T_1/I_1},
\end{equation}
and symmetrically at the cold end $j=n$ with $T_1\to T_n$; here $\gamma>0$ is the
coupling strength. The bath \eqref{eq:bath} is constructed so that at equilibrium
$T_1=T_n=T$ the system admits the explicit Gibbs invariant density
\begin{equation}\label{eq:gibbs}
  \rho_{\mathrm{eq}}(\bI,\bphi) \;\propto\; \exp\!\big(-H(\bI,\bphi)/2T\big).
\end{equation}

\begin{remark}[Choice of thermostat]\label{rem:bath}
The analysis of \cite{HLNS} uses a different thermostat---dissipation of the form
$\gamma(T_j-I_j^3)$ with a tailored multiplicative angle noise---chosen to make
the Lyapunov construction tractable; under that choice the Gibbs measure
\eqref{eq:gibbs} is \emph{not} invariant, even at equal temperatures. For our
numerical study we instead use the fluctuation--dissipation bath
\eqref{eq:bath}, whose equilibrium measure is the explicit Gibbs density
\eqref{eq:gibbs}; this provides a stringent check on the Fokker--Planck solver
(\cref{sec:fp-numerics}). The reduced-system analysis of the stabilization
mechanism in \cref{sec:stabilization} retains the original $\gamma(T-I^3)$ bath
of \cite{HLNS} for direct comparison with the proof.
\end{remark}

\subsection{Known results and open problems}\label{sec:known}

For the three-mode chain ($n=3$), in which only the middle mode $I_2$ is
unforced, Hani, Li, Nahmod and Staffilani \cite{HLNS} establish the existence,
uniqueness and ergodicity of the NESS: when the two bath temperatures are
sufficiently different, the system possesses a unique invariant probability
measure $\pi$, and the law of the process converges to $\pi$ in total variation
at a power-law rate in time. The proof rests on a novel construction of a
Lyapunov function through an elliptic Feynman--Kac equation. The central
difficulty is that, unlike microscopic heat conduction where energy simply flows
from hot to cold, here the transfer of energy to and from the unforced middle
mode depends sensitively on the phase variables being ``correct'': the system
must be prevented both from \emph{overheating} (energy trapped in $I_2$) and from
\emph{freezing} (energy blocked at $I_2$), and which of these threatens depends on
whether $I_2$ is large or small. This phase-dependent stabilization mechanism is
the analytical heart of \cite{HLNS} and the object of our \cref{sec:stabilization}.

Two limitations of the analytical theory motivate the present work. First, the
results of \cite{HLNS} are confined to $n=3$; the authors explicitly leave longer
chains open, noting that already the three-mode case is demanding. Second, the
NESS is not given in closed form, so quantitative properties---energy profiles,
transport coefficients, local statistics, the relaxation spectrum---are
inaccessible to the existence theory. Both gaps are naturally addressed
numerically.

\subsection{Contributions and organization}\label{sec:contributions}

This paper is a numerical study of the NESS of the energy cascade model
\eqref{eq:sde-bulk}. After collecting the thermodynamic and spectral framework in
\cref{sec:thermo,sec:fp,sec:methods}, we present two groups of results.

\Cref{sec:long-chain} treats the \emph{long chain}, the regime left open by
\cite{HLNS}. We characterize the nonequilibrium steady state---the action
profile and an approximate local equilibrium, in which adjacent-mode pair
marginals match a rescaled equilibrium marginal at the local kinetic temperature
(without strict local-Gibbs structure)---and we study the transport of the
conserved action current. The numerics indicate an anomalously fast decay of
the mean current with chain length, i.e. an insulating action-transport regime.

\Cref{sec:short-chain} returns to the \emph{three-mode chain} for a detailed
study. We solve the stationary Fokker--Planck equation by a mesh-free,
data-driven neural-network method \cite{ZhaiDobsonLi}, validating it against the
exact Gibbs measure \eqref{eq:gibbs} in equilibrium; we then examine the
nonequilibrium steady state, the phase-locking stabilization mechanism, and
trajectory-based estimates of the slow relaxation mode. The numerical tools---a
high-throughput SIMD Monte Carlo integrator, a kernel density estimator for the
reference measure, and a trajectory-based backward solver for the dominant
eigenpair---are described as they arise.


\section{Preliminary}

This section collects the thermodynamic observables we measure for the long
chain (\cref{sec:thermo}), the spectral (Fokker--Planck) framework for the
short chain (\cref{sec:fp}), the numerical methods used to solve the stationary
Fokker--Planck equation and to extract eigenfunctions of the generator
(\cref{sec:methods}), and the validation protocol for the Monte Carlo
transport computations (\cref{sec:mc-validation}).

\subsection{Nonequilibrium thermodynamics of the chain}\label{sec:thermo}

The Hamiltonian flow \eqref{eq:ham-flow} conserves the total mass
$M=\sum_k I_k$, and the evolution of each action obeys a discrete continuity
equation
\begin{equation}\label{eq:continuity}
  \dot I_j \;=\; J_j - J_{j+1},
  \qquad
  J_j \;:=\; 4\, I_{j-1} I_j \sin\theta_j ,
\end{equation}
so that $J_j$ is the \emph{action current} flowing from mode $j-1$ into mode
$j$ (with $J_1=J_{n+1}=0$ for the closed chain). In the nonequilibrium steady
state (NESS) the expected actions are stationary, $\E[\dot I_j]=0$, hence the
mean current is uniform along the chain, $\E[J_j]\equiv \E[J(n)]$ for all interior
$j$.

\begin{definition}[Thermal conductivity and Fourier's law]\label{de:fourier}
Assume $T_1>T_n$, so the positive action-current direction points from the hot
left bath toward the cold right bath. Let
\[
  \nabla T(n):=\frac{T_n-T_1}{n}<0
\]
denote the effective temperature gradient in that positive spatial direction.
We say, by analogy with Fourier's law, that the chain has finite action
conductivity $\kappa$ if
\begin{equation}\label{eq:fourier}
  \lim_{n \rightarrow \infty}\frac{\E[J(n)]}{\nabla T(n)} \;=\; -
  \kappa \in (0,\infty).
\end{equation}

Let $\kappa_n = -\E[J(n)]/\nabla T(n)$ be the finite-chain action
conductivity. Anomalous transport corresponds to $\kappa_n\sim n^{\alpha}$ with
$\alpha\neq0$:
$\alpha>0$ is superdiffusive and $\alpha<0$ (so $\kappa_n\to0$) signals an
insulator.
\end{definition}

When one measures the actual Hamiltonian energy exchanged with the baths, two
further observables characterize the NESS. The \emph{entropy production rate}
associated with the two heat baths is
\begin{equation}\label{eq:entropy-prod}
  \dot S \;=\; -\frac{\E[\dot Q_1]}{T_1} - \frac{\E[\dot Q_n]}{T_n} \;\ge\; 0,
\end{equation}
where $\dot Q_1,\dot Q_n$ are the mean heat fluxes delivered by the baths; in
the NESS the energy balance forces $\E[\dot Q_1]=-\E[\dot Q_n]=:\E[\dot Q]$ and
$\dot S = \E[\dot Q]\,(1/T_n - 1/T_1)$. The second law $\dot S\ge 0$ is refined
by the steady-state \emph{fluctuation theorem} of Gallavotti--Cohen
\cite{GallavottiCohen}: under its usual chaoticity and large-deviation
hypotheses, writing $\Sigma_\tau$ for the entropy produced over a window of
length $\tau$, the corresponding densities satisfy formally
\begin{equation}\label{eq:ft}
  \lim_{\tau\to\infty}\frac{1}{\tau}\,
  \ln\frac{p_\tau(a)}{p_\tau(-a)} \;=\; a .
\end{equation}
In this paper the transport measurements below concern the conserved
\emph{action} current. A full bath-energy and entropy-production measurement is
left as a separate numerical diagnostic and is not used as evidence for the
main claims.

\subsection{Fokker--Planck equation and spectral decomposition}\label{sec:fp}

Write $x=(\bI,\bphi)$ and let the system \eqref{eq:sde-bulk} have drift $f(x)$
and diffusion matrix $\Sigma(x)=\sigma(x)\sigma(x)^{\!\top}$. The probability
density $u(x,t)$ evolves by the (forward) Fokker--Planck, or Kolmogorov forward,
equation $\partial_t u = \Ladj u$, where
\begin{equation}\label{eq:fp-operator}
  \Ladj u = -\sum_i \partial_{x_i}\!\big(f_i\,u\big)
            + \tfrac12\sum_{i} \partial^{2}_{x_i}\!\big(\Sigma_{ii}\,u\big).
\end{equation}
For our model the diffusion is diagonal and degenerate: only the bath modes
carry noise, with
\begin{equation}\label{eq:diffusion}
  \tfrac12\Sigma_{I_j I_j} = 4\gamma T_j I_j, \qquad
  \tfrac12\Sigma_{\varphi_j\varphi_j} = \gamma T_j/I_j,
  \qquad j\in\{1,n\},
\end{equation}
and zero otherwise in the full $(\bI,\bphi)$ coordinates. In reduced
coordinates built from doubled bond angles, for instance
$\theta_2=2(\varphi_2-\varphi_1)$ at the left boundary, the angular diffusion
coefficient is multiplied by $2^2$, giving
$\tfrac12\Sigma_{\theta_2\theta_2}=4\gamma T_1/I_1$ (and similarly at the right
boundary).

\begin{remark}[Degeneracy / hypoellipticity]\label{rem:hypo}
Because noise enters only at $j\in\{1,n\}$, $\Sigma$ has rank $4$ in the
$2n$-dimensional state space; the interior modes are forced only through the
Hamiltonian coupling. Existence, uniqueness and ergodicity of the invariant
measure for the $n=3$ system are established in \cite{HLNS} via a
Feynman--Kac--Lyapunov argument together with H\"ormander's condition; the long
chain $n>3$ is left there as an open problem and is the subject of
\cref{sec:long-chain}.
\end{remark}

The stationary density solves $\Ladj u = 0$ with $\int u\dd x = 1$. The
long-time relaxation toward it is governed by the (backward) generator $\Lgen$,
the formal adjoint of $\Ladj$,
\begin{equation}\label{eq:generator}
  \Lgen g = \sum_i f_i\,\partial_{x_i} g
          + \tfrac12\sum_i \Sigma_{ii}\,\partial^2_{x_i} g,
  \qquad
  (\E_x[g(X_t)]) = (e^{t\Lgen} g)(x).
\end{equation}
Under ergodicity the spectrum of $\Lgen$ lies in $\{\mathrm{Re}\,\lambda\le 0\}$
with a simple eigenvalue $\lambda_0=0$ (constant eigenfunction). Ordering the
remaining eigenvalues by real part, the \emph{spectral gap}
$\lambda_R := \max\{\mathrm{Re}\,\lambda : \lambda\neq 0\}<0$ sets the slowest
relaxation rate, and for an observable $g$ orthogonal to constants
\begin{equation}\label{eq:decay}
  \E_x[g(X_t)] - \E_{\mathrm{eq}}[g]
  \;\sim\; Q(x)\,e^{\lambda_R t}\quad (t\to\infty),
\end{equation}
where $Q$ is the corresponding (sub)dominant eigenfunction, $\Lgen Q =
\lambda_R Q$. A purely real $\lambda_R$ indicates an overdamped relaxation; a
nonzero imaginary part $\lambda_I$ indicates damped oscillation with frequency
$\lambda_I$.

\subsection{Numerical methods}\label{sec:methods}

\subsubsection{A data-driven neural-network Fokker--Planck solver}
Classical grid-based PDE solvers are impractical for \eqref{eq:fp-operator} in
dimension $2n$: the domain is unbounded with no natural boundary condition, and
the cost grows exponentially with dimension. We follow the mesh-free,
data-driven solver of Zhai--Dobson--Li \cite{ZhaiDobsonLi} (building on
\cite{Li2019,DobsonLiZhai}), which represents the stationary density by a neural
network $\tilde u(x;\vartheta)$ trained against the unconstrained loss
\begin{equation}\label{eq:nn-loss}
  \mathcal{J}(\vartheta)
  = \underbrace{\frac{1}{N_X}\sum_{i}\big(\Ladj \tilde u(x_i;\vartheta)\big)^2}_{\text{PDE residual } L_1}
  + \underbrace{\frac{1}{N_Y}\sum_{j}\big(\tilde u(y_j;\vartheta)-v(y_j)\big)^2}_{\text{MC reference } L_2},
\end{equation}
where $\Ladj\tilde u$ is evaluated by automatic differentiation and $v$ is a
low-accuracy Monte Carlo (MC) estimate of the density at reference points
$\{y_j\}$. The operator term $L_1$ regularizes the network toward solutions of
$\Ladj u=0$ and sharply reduces the amount and accuracy of MC data required;
the method tolerates spatially uncorrelated noise in $v$ \cite{ZhaiDobsonLi}.

We adapt the solver to \eqref{eq:sde-bulk} in two ways. First, the angles are
periodic, so the network input is the \emph{periodic encoding}
$(\bI,\cos\theta_1,\sin\theta_1,\dots)$ rather than the raw angles. Second, the
diffusion \eqref{eq:diffusion} is degenerate (\cref{rem:hypo}); we therefore
include collocation points in the zero-density region and constrain them by the
operator $L_1$ alone (no fabricated reference data there). The reference density
$v$ is obtained by a high-throughput MC simulation of \eqref{eq:sde-bulk}
together with a kernel density estimate; see \cref{sec:fp-numerics}.

\subsubsection{Eigenfunctions of the generator}
To probe \eqref{eq:decay} we estimate the relaxation rate and slow-mode
amplitude visible to a chosen observable directly from trajectories. For a
localized observable $g$ (we use $g=\cos\theta_1$) we evolve a large ensemble
of trajectories from each of many initial points $x$ and record the conditional
expectation $\E_x[g(X_t)]$ as a function of $t$. Fitting a late-time
exponential window gives an
observable-dependent estimate of the dominant decay rate visible to $g$; if a
damped oscillatory fit is also stable, it gives an estimate of an imaginary
part. The fitted amplitude gives a pointwise sample $Q(x)$ of the associated
slow mode. A second network is then trained to represent $Q$ on the whole
domain, using a data loss against the sampled $Q(x)$ together with a PDE loss
enforcing the eigenvalue equation $\Lgen Q=\lambda_R Q$. These fits are
diagnostic rather than a proof that the full generator spectrum is real. Results
are reported in
\cref{sec:eigenfunction}.

\subsection{Monte Carlo validation and uncertainty protocol}\label{sec:mc-validation}

The long-chain transport claims use a separate canonical current accumulator
for the Gibbs-preserving SDE \eqref{eq:sde-bulk}. The accumulator is written in
double precision, uses deterministic per-trajectory random-number streams, and
records one independent time-averaged current per trajectory together with four
equal subwindows. We report standard errors across independent trajectory
averages and use trajectory-cluster bootstrap resampling for scaling-exponent
confidence intervals; the subwindows are used only as finite-time diagnostics,
not as independent trajectories.

Before production runs the code is subjected to four validation gates. First,
the optimized build is warning-clean and a sanitizer smoke run completes.
Second, repeating the same seeded run with one and two OpenMP threads gives
byte-identical samples, profiles, and burn-in traces. Third, an
equal-temperature run gives zero mean action current within confidence
intervals, while swapping the hot and cold baths reverses the current with the
same magnitude. Fourth, an equal-temperature SDE action profile agrees with an
independent Metropolis sampler of the Gibbs law \eqref{eq:gibbs}, with maximum
discrepancy $1.10$ combined standard errors in the validation run. These gates
check the sign convention, the fluctuation--dissipation normalization, and the
statistical independence assumptions used in the current estimates.

The timestep study is used only to choose a stable production step, not to hide
a continuum extrapolation. The coarse step $\Delta t=10^{-3}$ is acceptable at
$n=10$ but increasingly biased at larger $n$; at $n=40$ it underestimates the
current by $17.6\%$ relative to $\Delta t=5\times10^{-4}$. In contrast, the
$n=40$ comparison between $\Delta t=5\times10^{-4}$ and
$2.5\times10^{-4}$ differs by $3.9\%$, or $0.60$ pooled standard errors. The
reported scaling exponent is therefore a finite-size numerical law conditional
on the validated timestep and measurement window, not an asymptotic theorem.

\section{Statistical physics properties of the long cascade chain}\label{sec:long-chain}

The rigorous analysis of \cite{HLNS} treats the three-mode chain and lists the
longer chains $n>3$ as open. This section studies those long chains
numerically under the Gibbs-preserving two-bath dynamics \eqref{eq:bath}: the
structure of the nonequilibrium steady state (\cref{sec:ness-long}), the
transport of the conserved action current and the associated finite-chain
conductivity (\cref{sec:thermal-conductivity}), and finite-time current
fluctuations. We keep the terminology ``action current'' throughout, because
\eqref{eq:continuity} is a continuity equation for the conserved action
$M=\sum_j I_j$; Hamiltonian energy exchange with the baths is a distinct
observable.

\paragraph{Simulation.}
All long-chain data are produced by a single Euler--Maruyama integrator of
\eqref{eq:sde-bulk} that processes $W=16$ independent trajectories in lockstep
(one SIMD register per arithmetic operation), with the phases wrapped to
$[-\pi,\pi)$ at every step. Unless otherwise stated we use $\gamma=0.1$ and the
closed two-bath chain \eqref{eq:bath}. The local-equilibrium histograms use
long on-the-fly accumulation runs, while the current-scaling runs use an
independent canonical current accumulator with double precision, deterministic
trajectory seeds, four equal current blocks per trajectory, and explicit
timestep and burn-in diagnostics as described in \cref{sec:mc-validation}.

\subsection{Nonequilibrium steady state}\label{sec:ness-long}

\paragraph{Stabilization.}
We regard a run as stabilized when the sampled expectations (in particular the
action profile $\langle I_j\rangle$ and the bond-angle statistics) cease to
change with increasing sample size. Because the cascade transports energy
slowly, the relaxation time is long and the burn-in must be chosen with care:
a short run can report a spuriously uniform profile that has not yet developed
the steady gradient. We return to this point below.

\paragraph{Energy profile.}
\Cref{fig:profile} shows the mean action profiles $\langle I_j\rangle$ for
closed chains with $T_1=10$, $T_n=2$ at $n=25,50,100$. The profiles are
monotone after a short boundary layer at the hot contact and show the sharpest
drop at the cold bath. For $n=25$, the profile decreases from
$\langle I_1\rangle=0.738$ at the hot end to $\langle I_{13}\rangle=0.573$ near
mid-chain and $\langle I_{25}\rangle=0.165$ at the cold contact. The
corresponding endpoint values are $(0.514,0.106)$ at $n=50$ and
$(0.362,0.073)$ at $n=100$. These profile estimates are the data used below to
define the local kinetic temperature in the LTE comparison.

\begin{figure}[ht]
  \centering
  \includegraphics[width=0.82\textwidth]{action_profiles.pdf}
  \caption{Steady-state action profiles of the long closed chain. The
  horizontal coordinate is normalized as $j/(n-1)$ to compare
  $n=25,50,100$ on the same axis. The sample counts are
  $5.07\times10^8$, $1.01\times10^9$, and $4.06\times10^9$, respectively.}
  \label{fig:profile}
\end{figure}

\paragraph{Cascade geometry.}
The interpretation of the chain as an energy cascade rests on the frequency
embedding of \cite{CKSTT}: each generation of the resonant ``nuclear family''
construction places a group of modes on $\mathbb{Z}^2$ by exact rational
(Pythagorean) rotations, so that the abstract chain position $j$ corresponds to
a concrete set of frequencies $\Lambda_j$. \Cref{fig:embedding} plots the
locations of the first five generations used in the CKSTT construction. The
open-chain terminal-energy scaling is not used as evidence in the present
paper; establishing it would require a separate convergence study under the
same corrected thermostat used for the closed-chain transport measurements.

\begin{figure}[ht]
  \centering
  \includegraphics[width=0.62\textwidth]{cascade_embedding.pdf}
  \caption{Cascade geometry: the first five generations in the CKSTT resonant
  frequency embedding. Points are Fourier modes $K\in\mathbb{Z}^2$ (or exact
  rational intermediate coordinates in the displayed construction), and
  segments indicate resonant family transitions between adjacent generations.}
  \label{fig:embedding}
\end{figure}


\paragraph{Local thermodynamic equilibrium.}\label{par:lte}
A central question is whether, despite the global nonequilibrium drive,
adjacent modes in the bulk are described \emph{locally} by an equilibrium
measure. Let $q(I_a,I_{a+1},\theta_a)$ denote the joint pair density in the
NESS, with $(I_a,I_b)=(I_j,I_{j+1})$, and let
$p_T(I_a,I_{a+1},\theta_a)$ be the corresponding marginal of the equilibrium
measure at temperature $T$, whose full density is proportional to
$\exp[-H/(2T)]$ under the convention of \cref{rem:convention} (integrating out
all other modes). We accumulate $q$ and $p_T$ on
$80^3$ histograms over $[0,4]^2\times[-\pi,\pi)$ from $\sim10^9$ samples per
site, at interior pairs of the closed chain ($T_1=10,T_n=2$) for
$n=25,50,100$.

First, neither $q$ nor $p_T$ equals the two-mode Gibbs density
$\exp\!\big[-(\tfrac12(I_a{+}I_b)^2-\tfrac14(I_a^2{+}I_b^2)+I_aI_b\cos\theta)\big]$:
comparing $-\log(\text{PDF})$ to this form, the NESS and the equilibrium
marginal deviate in the \emph{same} way (both rise along the axes where the
two-mode form is flat, and their $\theta$-modulation is weaker). Since the
equilibrium marginal already deviates identically, the discrepancy comes from
marginalizing the nonlocal term $\tfrac12(\sum_j I_j)^2$, not from the
nonequilibrium drive: there is no strict local Gibbs structure, even at
equilibrium.

Second, the NESS marginal is nonetheless well approximated by a
\emph{rescaled} equilibrium marginal. We fit $\log q = x\log p_6 + c$ by
weighted least squares, weighting each bin by its count (the inverse variance
of the log-density for multinomial counts); statistical errors on $x$ are
below $3\times10^{-4}$ and are negligible against the systematic spread
discussed below. Restricted to the core ($I_a,I_b<2.5$) the fit is excellent,
with weighted $R^2$ from $0.986$ to $0.999$ (\cref{tab:lte}), and the rescaled
$p_6^{\,x}$ matches $q$ closely (\cref{fig:lte-resid}).

We caution against reading the slope as a temperature through the naive
relation $6/x$. A control fit between two simulated equilibria,
$\log p_{7.12}=x\log p_6+c$, returns $x=0.926$ rather than the naive
$6/7.12=0.843$: the marginals do not scale as $p_T=p_6^{\,6/T}$,
because the mean field $\langle\sum_j I_j\rangle\propto\sqrt{nT}$ shifts with
temperature. Hence $6/x$ systematically underestimates the effective
temperature.

Third, the effective temperature is the local kinetic temperature. Define
\[
  T_{\mathrm{kin}}(a)=
  6\,\frac{\langle I_a\rangle_{\mathrm{NESS}}}
          {\langle I_a\rangle_{\mathrm{eq}}}
\]
directly from the action profiles, independent of the marginal fit. Matching
$q$ to an actual simulated equilibrium confirms
this scale: at $a/n=0.48$ the NESS marginal fitted against the true $p_{7.12}$
gives $x=1.01$, i.e.\ $q\approx p_{7.12}$, within a few percent of
$T_{\mathrm{kin}}=6.85$ at that site (and consistent to $0.2\%$ with the
$q$--$p_6$ and $p_{7.12}$--$p_6$ fits combined). The kinetic temperature, the
equilibrium-matching temperature, and the naive $6/x$ agree to within
$5$--$10\%$, the spread reflecting the residual departure of $q$ from a true
equilibrium marginal. To leading order, then, the dependence of the effective
temperature on $(T_L,T_R,a/n)$ is simply the shape of the temperature profile;
across boundary pairs of equal mean ($(8,4),(9,3),(10,2)$) it varies
monotonically with the imposed gradient.

\begin{table}[ht]\centering
\begin{tabular}{ccccc}
\hline
$a/n$ & $n$ & $x$ & $R^2_{\mathrm w}$ & $6/x$ \quad $T_{\mathrm{kin}}$ \\
\hline
       & $25$  & $0.8203$ & $0.9970$ & $7.31\quad 7.68$ \\
$0.24$ & $50$  & $0.8269$ & $0.9984$ & $7.26\quad 7.79$ \\
       & $100$ & $0.8429$ & $0.9967$ & $7.12\quad 7.82$ \\
\hline
       & $25$  & $0.9360$ & $0.9943$ & $6.41\quad 6.85$ \\
$0.48$ & $50$  & $0.9308$ & $0.9988$ & $6.45\quad 7.03$ \\
       & $100$ & $0.9304$ & $0.9989$ & $6.45\quad 7.01$ \\
\hline
       & $25$  & $1.0933$ & $0.9862$ & $5.49\quad 5.80$ \\
$0.72$ & $50$  & $1.0610$ & $0.9919$ & $5.65\quad 6.04$ \\
       & $100$ & $1.0412$ & $0.9980$ & $5.76\quad 6.02$ \\
\hline
\multicolumn{5}{l}{\emph{Controls} ($n=25$):}\\
\multicolumn{2}{l}{$p_{7.12}$ vs.\ $p_6$ \ (ceiling)}      & $0.9261$ & $0.9997$ & --- \\
\multicolumn{2}{l}{NESS vs.\ $p_{7.12}$ \ ($a/n{=}0.48$)}  & $1.0093$ & $0.9944$ & --- \\
\hline
\end{tabular}
\caption{Weighted least-squares fit $\log q=x\log p_6+c$ over the core
($I_a,I_b<2.5$): slope $x$, weighted $R^2$, the naive temperature $6/x$, and
the independently measured kinetic temperature $T_{\mathrm{kin}}$, at fixed
$a/n$ for $n=25,50,100$. Statistical errors on $x$ are below $3\times10^{-4}$,
negligible against the $5$--$10\%$ spread among temperature estimators. The
two control fits (two equilibria, giving the ceiling; and NESS against the
matched equilibrium $p_{7.12}$, $x\approx1$) are listed below. The naive
$6/x$ lies systematically below $T_{\mathrm{kin}}$ because the marginals do not
follow the naive Gibbs scaling.}
\label{tab:lte}
\end{table}

Fourth, the approximation improves as the chain lengthens. The control fit
between two equilibria ($p_{7.12}$ vs.\ $p_6$) gives a weighted $R^2=0.9997$,
the ceiling attainable within the equilibrium family. At fixed $a/n$ the core
$R^2$ rises toward this ceiling with $n$---most clearly at the cold,
large-gradient sites ($a/n=0.72$: $0.986\to0.992\to0.998$ for $n=25,50,100$)
and at mid-chain ($a/n=0.48$: $0.994\to0.999\to0.999$), while the hot site
($a/n=0.24$) is already near the ceiling at $n=25$. A small gap persists at
all $n$ (about $0.001$--$0.002$ at $n=100$); in this sense $q\to p_{\beta_0}$
in the long-chain interior without becoming exactly equal to it.

That residual gap is the genuine deviation from local equilibrium. It has two
parts: a curvature---core and tail carry slightly different effective
temperatures, so a single $x$ cannot fit both---and a genuine nonequilibrium
component, an odd-in-$\theta$ part of $q$ that no equilibrium density (which
depends on $\theta$ only through $\cos\theta$) can reproduce, tied to the
action flux of \cref{sec:thermal-conductivity}. Neither is removed by
rescaling to a single temperature, so strict local Gibbs structure is not
restored. In summary, in the long-chain interior the NESS pair marginal is
well described by an equilibrium marginal at the local kinetic temperature,
increasingly so as $n$ grows, but strict local Gibbs structure does not hold
exactly in either ensemble.

\begin{figure}[ht]
  \centering
  \includegraphics[width=0.82\textwidth]{lte_residual_midchain.pdf}
  \caption{Residual of the rescaled-equilibrium fit,
  $\log q-(x\log p_6+c)$, for the mid-chain pair $a/n=0.48$ at fixed
  $\theta\approx0$. Low-count bins (fewer than $200$ samples in either the NESS
  or equilibrium histogram on the displayed slice) are masked. The fitted slopes
  and weighted $R^2$ values agree with the core values in \cref{tab:lte}; the
  remaining structured residual is the finite, non-Gibbs component discussed in
  the text.}
  \label{fig:lte-resid}
\end{figure}

For a more geometric view of the same approximation error,
\cref{fig:lte-resid-mesh} shows the residual surface on the displayed
$\theta$-slice using the MATLAB \texttt{compare\_residual.m} masking convention.
This diagnostic includes the short $n=15$ chain and two longer chains at a
comparable interior location. The surfaces show a coherent, mostly
low-curvature core together with a sharper boundary/tail ridge: the dominant
effect is therefore a smooth mismatch between the NESS marginal and a single
rescaled equilibrium marginal, rather than unstructured Monte Carlo noise. The
figure is used as a structural residual diagnostic; the numerical slopes in
\cref{tab:lte} are computed by the fixed weighted-core estimator described
above.

\begin{figure}[ht]
  \centering
  \includegraphics[width=\textwidth]{report_assets/compare_residual_mesh.pdf}
  \caption{Three-dimensional residual surfaces
  $\log q-(x\log p+c)$ for the pair-marginal local-equilibrium comparison,
  reproduced from the MATLAB \texttt{compare\_residual.m} convention. The
  displayed slice uses $80^3$ histograms, the mask $(Q>50)$ and $(P>50)$ for
  the fit, and an interior hot-side pair for $n=15,25,50$. The short-chain
  $n=15$ residual has the same qualitative shape as the longer chains,
  supporting the interpretation that the residual is a structured finite
  nonequilibrium correction rather than a plotting artifact.}
  \label{fig:lte-resid-mesh}
\end{figure}

\subsection{Thermal conductivity}\label{sec:thermal-conductivity}

In the closed-chain NESS the mean action current $\E[J]$ of
\eqref{eq:continuity} is uniform along the chain. We measure how it scales with
the chain length $n$ using the Gibbs-preserving bath \eqref{eq:bath}. The
production runs use $T_1=10$, $T_n=2$, $\gamma=0.1$, timestep
$\Delta t=5\times10^{-4}$, and $1024$ independent trajectories for each chain
length; the burn-in is scaled with the observed relaxation time
($1000,1280,2880,5120$ for $n=10,20,30,40$). Each trajectory is averaged over a
window of length $200$ after burn-in. The measured currents are
\[
\begin{array}{c|cccc}
n & 10 & 20 & 30 & 40\\ \hline
\E[J(n)] &
0.39252 & 0.11917 & 0.05457 & 0.02975\\
\mathrm{SE} &
0.00189 & 0.00092 & 0.00062 & 0.00048
\end{array}
\]
where the standard errors are computed across independent trajectory averages.
A log--log fit gives
\begin{equation}\label{eq:flux-scaling}
  \E[J(n)] \;=\; 28.75\,n^{-1.850},
  \qquad R^2=0.9980,
\end{equation}
with a trajectory-bootstrap $95\%$ confidence interval
$[-1.870,-1.830]$ for the exponent. Thus $\E[J(n)]\,n\to0$: the current decays
faster than the Fourier scaling $1/n$. By \cref{de:fourier},
\[
  \kappa_n=\frac{n\,\E[J(n)]}{T_1-T_n}
  \sim 3.59\,n^{-0.850},
\]
so the long chain behaves as an action-current insulator. This is the opposite
of the divergent conductivity familiar from FPU-type lattices
\cite{LepriLiviPoliti}, although the transported conserved quantity here is the
action rather than the Hamiltonian energy.

Two numerical checks accompany \eqref{eq:flux-scaling}. First, every production
trajectory is split into two equal measurement windows; the largest
first-half/second-half discrepancy among the four lengths is only $0.90$
paired standard errors. Second, a timestep pilot shows that the coarse
$\Delta t=10^{-3}$ run underestimates the $n=40$ current, whereas
$5\times10^{-4}$ and $2.5\times10^{-4}$ agree within the Monte Carlo error
(relative difference about $3.9\%$, $0.60$ pooled standard errors, in the
64-trajectory pilot). The production exponent is therefore based on
$\Delta t=5\times10^{-4}$.

\begin{figure}[ht]
  \centering
  \includegraphics[width=0.72\textwidth]{experiments/flux_validation/production_dt5e-4/flux_primary_scaling.pdf}
  \caption{Anomalous scaling of the action current. Mean finite-time averaged
  current $\E[J(n)]$ vs.\ chain length $n$ (log--log) for the corrected
  Gibbs-preserving SDE. Error bars are normal $95\%$ confidence intervals for
  the trajectory mean; the fitted exponent is $-1.850$ with bootstrap
  confidence interval $[-1.870,-1.830]$.}
  \label{fig:flux-scaling}
\end{figure}

\begin{table}[ht]\centering
\small
\begin{tabular}{cccccc}
\toprule
$n$ & $\Delta t_{\rm fine}$ & $\Delta t_{\rm coarse}$ &
fine mean & coarse mean & coarse--fine / fine \\
\midrule
10 & $5.0{\times}10^{-4}$ & $10^{-3}$ &
0.39252 & 0.39059 & $-0.5\%$ \\
20 & $2.5{\times}10^{-4}$ & $5.0{\times}10^{-4}$ &
0.11815 & 0.11917 & $+0.9\%$ \\
20 & $5.0{\times}10^{-4}$ & $10^{-3}$ &
0.11917 & 0.11529 & $-3.3\%$ \\
30 & $5.0{\times}10^{-4}$ & $10^{-3}$ &
0.05457 & 0.04943 & $-9.4\%$ \\
40 & $2.5{\times}10^{-4}$ & $5.0{\times}10^{-4}$ &
0.02863 & 0.02975 & $+3.9\%$ \\
40 & $5.0{\times}10^{-4}$ & $10^{-3}$ &
0.02975 & 0.02450 & $-17.6\%$ \\
\bottomrule
\end{tabular}
\caption{Timestep sensitivity for the current estimator. The production step
$\Delta t=5\times10^{-4}$ is retained because the coarser
$10^{-3}$ step visibly biases the longer chains, whereas the available
finer-step pilots at $n=20$ and $n=40$ agree with the production step within
Monte Carlo error. The table reports pilot means, so it is a convergence
diagnostic rather than an additional scaling fit.}
\label{tab:dt-sensitivity}
\end{table}

Beyond the mean, we examine the distribution of finite-time averaged currents
\[
  \overline J_\tau=\frac1{\tau}\int_{t_0}^{t_0+\tau}J(t)\,\dd t .
\]
The production runs store four disjoint current blocks per trajectory, giving
non-overlapping windows $\tau=50,100,200$ from the same data. We report these
finite-window distributions descriptively (variance scaling, skewness,
kurtosis, and sign probability) and do not infer an asymptotic large-deviation
tail from a single window. At $n=40$, for example, the sign probability
$\Pr(\overline J_\tau<0)$ decreases from $0.233$ at $\tau=50$ to $0.026$ at
$\tau=200$, while the standardized skewness stays small ($0.010$, $0.024$,
$0.075$ for the three windows). The product
$\tau\,\mathrm{Var}(\overline J_\tau)$ decreases over these finite windows
($0.095$, $0.067$, $0.048$), indicating that $\tau=200$ is not yet an
asymptotic regime in which a single large-deviation rate can be read off
reliably.

\begin{figure}[ht]
  \centering
  \begin{subfigure}{0.48\textwidth}
    \centering
    \includegraphics[width=\textwidth]{experiments/flux_validation/production_dt5e-4/current_windows_n40_windows.pdf}
    \caption{$n=40$ standardized current windows.}
  \end{subfigure}\hfill
  \begin{subfigure}{0.48\textwidth}
    \centering
    \includegraphics[width=\textwidth]{experiments/flux_validation/production_dt5e-4/current_windows_variance_scaling.pdf}
    \caption{Finite-window variance scaling.}
  \end{subfigure}
  \caption{Finite-time current diagnostics. Left: standardized
  $\overline J_\tau$ at $n=40$ for $\tau=50,100,200$ compared with a standard
  normal reference. Right: $\tau\,\mathrm{Var}(\overline J_\tau)$ over the
  available windows. The finite-window variance has not plateaued, so the data
  are interpreted descriptively and are not used to infer an asymptotic
  current large-deviation rate.}
  \label{fig:current-windows}
\end{figure}

\subsection{Bath-energy entropy production: scope of the present numerics}
\label{sec:entropy}

The two-bath SDE also permits a bath-energy entropy-production diagnostic based
on \eqref{eq:entropy-prod}. That quantity requires accumulating the stochastic
Hamiltonian-energy exchange with each boundary bath, not merely the conserved
action current \eqref{eq:continuity}. We therefore treat the Gallavotti--Cohen
symmetry \eqref{eq:ft} as a future numerical check rather than a result of this
paper. The current data below establish the action-transport scaling and the
finite-time statistics of $\overline J_\tau$ only.


\section{Detailed property of the short chain}\label{sec:short-chain}

We now turn to the smallest nontrivial cascade, the three-mode system $n=3$,
for which existence, uniqueness and ergodicity of the NESS are rigorously known
\cite{HLNS}. The Hamiltonian and bath of \cref{sec:long-chain} depend on the
phases only through $\theta_1=2(\varphi_1-\varphi_2)$ and
$\theta_3=2(\varphi_3-\varphi_2)$, so the dynamics reduces from the six
variables $(\bI,\bphi)$ to the five-dimensional state
$x=(I_1,I_2,I_3,\theta_1,\theta_3)$, in which the middle mode $I_2$ carries no
bath. This low dimension lets us solve the stationary Fokker--Planck equation
directly and interrogate the invariant measure (\cref{sec:fp-numerics}), the
mechanism that stabilizes the NESS (\cref{sec:stabilization}), and the spectrum
of the generator (\cref{sec:eigenfunction}).

\subsection{Numerical solution of the nonequilibrium steady state}\label{sec:fp-numerics}

In the reduced coordinates the stationary density solves $\Ladj u=0$, with
$\Ladj$ as in \eqref{eq:fp-operator}: the drift is the $n=3$ specialization of
\eqref{eq:sde-bulk}, and the (degenerate) diffusion \eqref{eq:diffusion} acts
only on $(I_1,\theta_1)$ and $(I_3,\theta_3)$ with coefficients
$4\gamma T_1 I_1,\ 4\gamma T_1/I_1$ and $4\gamma T_3 I_3,\ 4\gamma T_3/I_3$. We
solve it with the neural-network method of \cref{sec:methods}. The reference
density $v$ is produced by a long SIMD Monte Carlo integration of
\eqref{eq:sde-bulk} followed by a kernel density estimate on a grid of boxes in
$x$; boxes in which the Monte Carlo records no sample are retained as
zero-density collocation points and are constrained by the operator term $L_1$
alone. The network outputs the log-density, $u(x)$ with $\rho=\exp(\sigma_0
u+\mu_0)$ for fixed normalization constants $(\mu_0,\sigma_0)$, takes the
periodic encoding $(I_1,I_2,I_3,\cos\theta_1,\sin\theta_1,\cos\theta_3,
\sin\theta_3)$ as input, and is trained in two phases: a data-only pre-training
against $\{v(y_j)\}$, followed by joint training in which the weight of the PDE
residual $L_1$ is ramped on gradually.

\paragraph{Validation against the equilibrium Gibbs measure.}
When $T_1=T_3=T$ the exact invariant density is the Gibbs measure
\eqref{eq:gibbs}, which provides a stringent test. Training at $T=5$ and
comparing the network density to \eqref{eq:gibbs} on the slices $I_1=I_2=I_3=I$
(\cref{fig:eq-validation}), the mean relative error is $1.5\%$, $2.6\%$ and
$6.4\%$ at $I=0.5,1,2$ respectively, with the pointwise $\log$-ratio
$\log(\rho_{\mathrm{NN}}/\rho_{\mathrm{Gibbs}})$ confined to a narrow band. The
solver thus recovers the known equilibrium measure across the range of actions
that carry appreciable probability.

\begin{figure}[ht]
  \centering
  \includegraphics[width=0.95\textwidth]{eq_validation.png}
  \caption{Equilibrium validation of the neural-network Fokker--Planck solver.
  At $T_1=T_3=5$ and $I_1=I_2=I_3=0.5$, the network density
  $\rho(\theta_1,\theta_3)$ (left) is compared with the analytic Gibbs measure
  (centre-left); the relative error (centre-right) and $\log$-ratio (right)
  stay within a few percent across the cell.}
  \label{fig:eq-validation}
\end{figure}

\paragraph{Nonequilibrium steady state.}
With distinct bath temperatures $T_1=2$, $T_3=8$ the invariant measure is no
longer of Gibbs form. \Cref{fig:neq-density} shows the network density on the
angular plane $(\theta_1,\theta_3)$ at several values of the common action: the
distribution concentrates around a phase-locked configuration and is visibly
distorted relative to the equilibrium measure. The departure from equilibrium
is quantified in \cref{sec:stabilization}.

\begin{figure}[ht]
  \centering
  \includegraphics[width=0.95\textwidth]{neq_density.png}
  \caption{Nonequilibrium steady-state angular density of the three-mode chain.
  Top row: network density $\rho(\theta_1,\theta_3)$ under the gradient
  $T_1=2,\,T_3=8$ at $I=0.3,0.5,1.0$; bottom row: the equilibrium
  ($T_1=T_3=5$) density at the same $I$. The nonequilibrium density is tilted
  and shifted off-centre, breaking the $\theta_1\leftrightarrow\theta_3$
  symmetry of the equilibrium case.}
  \label{fig:neq-density}
\end{figure}

\subsection{Mechanism of stabilization}\label{sec:stabilization}

The NESS is stabilized by a phase-locking mechanism that operates differently in
the high- and low-energy regimes \cite{HLNS}. At high energy the bond angles are
driven to a stable equilibrium near $2\pi/3$, escaping the unstable
configuration at $4\pi/3$; at low energy the angle $\theta_3$ remains near
$4\pi/3$, which prevents the middle action $I_2$ from collapsing to zero and
keeps the cascade alive. The reduced four-dimensional subsystem
$(I_1,I_3,\theta_1,\theta_3)$, obtained in the low-energy limit and equipped
with the $\gamma(T-I^3)$ thermostat of \cite{HLNS}, isolates this mechanism; we
solve its stationary Fokker--Planck equation by the same network method and
recover the predicted concentration of $\theta_3$.

We retain two source-traced diagnostics of the nonequilibrium drive at
$T_1=2,T_3=8$: symmetry breaking in the angular density and the phase-locking
branch predicted by the reduced equation.

\emph{(i) Symmetry breaking.} At equilibrium the density is symmetric under
$\theta_1\leftrightarrow\theta_3$. Out of equilibrium this symmetry is broken:
the relative asymmetry
$|\rho(\theta_1,\theta_3)-\rho(\theta_3,\theta_1)|/
(\rho(\theta_1,\theta_3)+\rho(\theta_3,\theta_1))$ is visibly nonzero on the
high-density part of the angular plane (\cref{fig:symmetry}). We report the
qualitative effect here and reserve the precise max/mean asymmetry for the
source-traced grid recomputation, since it depends sensitively on the density
mask and normalization.

\emph{(ii) Phase-locking angle.} The most probable bond angle is shifted from
its equilibrium value to
\begin{equation}\label{eq:phase-lock}
  \theta_0 \;=\; \pi-\arcsin\!\frac{\sqrt{4\gamma^2+3}-\gamma}{2(\gamma^2+1)},
\end{equation}
the stable branch that tends to $2\pi/3$ as $\gamma\to0$. We use the density
plots qualitatively to display concentration near the stable sector; no
quantitative peak-location table is used as an independent result here.

\begin{remark}\label{rem:eq47}
Equivalently, writing $s=\sin\theta_0$, the reduced angular fixed-point
condition gives the quadratic
\[
  4(1+\gamma^2)s^2+4\gamma s-3=0 .
\]
The positive root is
\[
  s=\frac{\sqrt{4\gamma^2+3}-\gamma}{2(1+\gamma^2)} .
\]
The inverse sine has two branches. The branch in \eqref{eq:phase-lock} lies in
$(\pi/2,\pi)$ and tends to $2\pi/3$ as $\gamma\to0$, matching the stable
high-energy phase-locking equilibrium; the complementary branch tends to
$\pi/3$ and is not the observed stable lock.
\end{remark}

Finally, current conservation through the unforced middle mode gives the
balance condition
$\E[\,I_2 I_1\sin\theta_1\,]+\E[\,I_2 I_3\sin\theta_3\,]\approx0$ in a
stationary state. We use this identity as a rerun diagnostic for the solver,
but do not report it as a separate quantitative result here.

\begin{figure}[ht]
  \centering
  \includegraphics[width=0.95\textwidth]{symmetry_breaking.png}
  \caption{$\theta_1\leftrightarrow\theta_3$ symmetry breaking in the
  nonequilibrium steady state. At $I_1=I_2=I_3=2$ and $T_1=2,\,T_3=8$, the
  network density $\rho(\theta_1,\theta_3)$ (left) and its transpose (centre)
  differ; the relative asymmetry $(\rho-\rho^{\mathsf T})/(\rho+\rho^{\mathsf
  T})$ (right) is shown with the colour scale clipped at $\pm0.15$ for
  readability. Since unmasked max/mean asymmetry is sensitive to low-density
  regions and network normalization, this panel is used as qualitative evidence
  of symmetry breaking rather than as a standalone quantitative estimator.}
  \label{fig:symmetry}
\end{figure}

\subsection{Eigenfunction of the generator}\label{sec:eigenfunction}

To probe the relaxation spectrum we apply the trajectory method of
\cref{sec:methods} with the localized observable $g=\cos\theta_1$. From each of
$4096$ initial points we evolve an ensemble and record
$\E_x[\cos\theta_1(X_t)]$; fitting the decay \eqref{eq:decay} returns, at each
initial point, an estimate of the observable-dependent relaxation rate and
amplitude.

At $\gamma=0.1$, pure-exponential fits
\[
  \E_x[\cos\theta_1(X_t)] \approx A(x)e^{\lambda t}+c(x)
\]
over the full recorded window $t\in[0,5]$ retain $2895$ of $4096$ curves under
the filter $-5<\lambda<0$, with median
\begin{equation}\label{eq:lambda_diag}
  \lambda_{\mathrm{diag}} \;\approx\; -0.934 .
\end{equation}
This number is reproducible from the saved backward-Monte-Carlo curves, but it
is not a window-independent spectral gap estimate: the retained median shifts
from $-1.60$ on $t\in[0,2]$ to $-0.68$ on $t\in[0.5,5]$
(\cref{fig:eigenvalues}). We therefore use $\lambda_{\mathrm{diag}}$ as a
diagnostic relaxation rate for the observable $\cos\theta_1$, not as a resolved
statement about the full generator spectrum.

The fitted amplitudes provide pointwise samples $Q(x)$ of the corresponding
slow-mode surrogate. A second network is trained to represent $Q$ subject to
the eigenvalue equation $\Lgen Q=\lambda_{\mathrm{diag}} Q$, using the generator
$\Lgen$ of \eqref{eq:generator} (without the divergence terms that distinguish
$\Ladj$). The reconstruction is only moderately accurate: the data
root-mean-square error is $0.38$ (relative error $0.71$) and the median
normalized PDE residual $|\Lgen Q-\lambda_{\mathrm{diag}} Q|/|Q|$ is $0.35$. We
report these candidly: the fitted amplitude data contain both signs and are
reconstructed from intrinsically noisy backward Monte Carlo data, making this
the hardest object computed in this work; the agreement is nonetheless
sufficient to exhibit the qualitative structure of $Q$ on representative
slices (\cref{fig:eigenfunction}).

\begin{figure}[ht]
  \centering
  \includegraphics[width=0.95\textwidth]{eigenvalue_scatter.png}
  \caption{Backward-Monte-Carlo relaxation-rate diagnostic at $\gamma=0.1$,
  using the observable $\cos\theta_1=\cos 2(\varphi_1-\varphi_2)$. Left:
  retained pure-exponential rates from $4096$ initial conditions over
  $t\in[0,5]$; $2895$ fits satisfy $-5<\lambda<0$, with median
  $\lambda_{\mathrm{diag}}\approx-0.934$. Right: median and interquartile range
  across several fitting windows, showing the window sensitivity that motivates
  treating this quantity as a diagnostic rate rather than as a high-accuracy
  spectral gap.}
  \label{fig:eigenvalues}
\end{figure}

\begin{figure}[ht]
  \centering
  \includegraphics[width=0.95\textwidth]{Q1_slices.png}
  \caption{Generator slow-mode surrogate. Network
  reconstruction of $Q(\theta_1,\theta_3)$ at $I_1=I_2=I_3=1,2,4$ (fitted
  $\lambda_{\mathrm{diag}}\approx-0.934$). The fitted amplitude data contain
  both signs; the displayed slices show the positive lobe concentrated near
  $\theta_1=\theta_3=0$. As noted in the text, the surrogate is reconstructed
  from noisy backward Monte Carlo data and is shown to exhibit the qualitative
  structure of $Q$ rather than as a high-accuracy spectral computation.}
  \label{fig:eigenfunction}
\end{figure}

\section{Conclusion and limitations}\label{sec:conclusion}

The numerical evidence supports three conclusions. First, the
Gibbs-preserving bath \eqref{eq:bath} gives a well-controlled statistical model
for the resonant NLS cascade chain: at equal bath temperatures the invariant
law agrees with the Gibbs density \eqref{eq:gibbs}, and at unequal temperatures
the same code produces a stationary action current with reproducible
time-reversal, timestep, and stationarity checks. Second, the long closed chain
does not exhibit Fourier action transport over the accessible lengths. The
corrected production runs give $\E[J(n)]\sim n^{-1.850}$, so the finite-chain
conductivity decays as $\kappa_n\sim n^{-0.850}$. The result should be read as
a finite-size numerical law for the conserved action current, not as a theorem
for Hamiltonian heat conduction. Third, the long-chain NESS has an approximate
local-equilibrium structure: adjacent-mode marginals are close to rescaled
equilibrium marginals at the local kinetic temperature, while the residual
odd-in-angle component and the equilibrium marginalization effect rule out
strict local Gibbs structure.

The short-chain computations complement this transport picture by resolving the
five-dimensional stationary density directly. The neural Fokker--Planck solver
recovers the known Gibbs equilibrium and, out of equilibrium, visualizes the
symmetry breaking and phase locking that underlie the stabilization mechanism
proved for the three-mode model in \cite{HLNS}. The generator-eigenfunction
calculation is less mature: it gives a consistent qualitative slow mode for the
observable $\cos\theta_1$, but the fitted eigenfunction has substantial
backward-Monte-Carlo noise and should be used as a diagnostic rather than as a
high-accuracy spectral computation.

Several limitations remain. The long-chain ergodicity theory is open, so the
stationarity of the long simulations is assessed numerically rather than
proved. The transport exponent is fitted over $n=10,20,30,40$ and should be
rechecked at larger lengths and smaller timesteps when computational resources
permit. The current statistics reported here concern finite-time averages of
the action current only; bath-energy entropy production and a genuine
Gallavotti--Cohen large-deviation test require a separate boundary-energy
accumulator. Finally, the open-chain terminal-energy scaling suggested by the
cascade geometry is not claimed here, because it has not yet been recomputed
with the corrected Gibbs-preserving workflow.

\section*{Numerical reproducibility summary}

\Cref{tab:repro-summary} summarizes the numerical audit trail used in this
revision. The table is intentionally limited to quantities that are either
regenerated from source files or checked against frozen validation artifacts.
It is not a substitute for the limitations above: the transport exponent is a
finite-size numerical fit, the short-chain neural networks are rerun from saved
models rather than retrained, and author declarations still require author
confirmation before submission.

\begin{table}[ht]\centering
\small
\setlength{\tabcolsep}{2pt}
\begin{tabular}{>{\raggedright\arraybackslash}p{0.18\textwidth}
                >{\raggedright\arraybackslash}p{0.24\textwidth}
                >{\raggedright\arraybackslash}p{0.31\textwidth}
                >{\raggedright\arraybackslash}p{0.17\textwidth}}
\toprule
Result family & Evidence bundle & Automated or source-traced check & Scope retained in manuscript \\
\midrule
Action-current scaling &
canonical simulator, production manifest, validation report, current-window
statistics &
frozen source/binary hashes; equal-temperature and hot/cold reversal gates;
trajectory-bootstrap exponent CI; split-window stationarity and timestep pilot &
finite-size action-current law over $n=10,20,30,40$ \\
\midrule
Long-chain profiles and local equilibrium &
figure metrics, histogram/profile hashes, LTE source-trace export &
profile sample counts; figure-source hashes; all entries of \cref{tab:lte}
recomputed from archived histograms and profiles &
local pair-marginal diagnostic, not strict local Gibbs convergence theorem \\
\midrule
Short-chain Fokker--Planck density &
archived notebook outputs, saved Keras models, TensorFlow rerun metrics &
equilibrium Gibbs-slice errors; qualitative symmetry breaking; phase-locking
branch; current-balance rerun diagnostic &
saved-model inference reproducibility, not full neural-network retraining \\
\midrule
Eigenfunction diagnostic &
backward-Monte-Carlo data, fit-window sensitivity JSON, saved-model rerun &
$0$--$5$ and alternate-window decay fits; eigen-surrogate data-fit RMSE;
reported PDE residual &
observable-dependent slow-mode diagnostic, not resolved spectral gap \\
\midrule
Manuscript-level claim audit &
claim-audit script plus JSON/Markdown registry &
registered core numerical/data claims checked against the artifacts above
before compilation &
local audit to rerun after final author and journal-format edits \\
\bottomrule
\end{tabular}
\caption{Compact reproducibility map for the numerical claims in the revised
manuscript. Exact file paths, hashes, and machine-readable values are recorded
in the artifacts listed in the data and code availability statement.}
\label{tab:repro-summary}
\end{table}

\section*{Data and code availability}

The code and numerical artifacts used for this revision are kept with the
project repository. Unless stated otherwise, paths below are relative to the
repository root. The key reproducibility artifacts are:
\begingroup
\small
\raggedright
\begin{itemize}[leftmargin=*]
  \item canonical current simulator: \path{flux/NLS_flux_canonical.cpp};
  \item current-analysis scripts:
  \path{flux/analyze_canonical_flux.py},
  \path{flux/analyze_current_windows.py},
  \path{flux/gibbs_mcmc_reference.py}, and
  \path{flux/compare_gibbs_sde.py};
  \item production current data and validation reports:
  \path{Paper/revision_2026-06-19/experiments/flux_validation/};
  \item figure-generation script for
  \cref{fig:profile,fig:embedding,fig:lte-resid}:
  \path{Paper/revision_2026-06-19/scripts/generate_manuscript_figures.py};
  \item source metrics for those figures:
  \path{Paper/revision_2026-06-19/manuscript_figure_metrics.json};
  \item source-trace export for the LTE table and short-chain diagnostics:
  \path{Paper/revision_2026-06-19/scripts/export_source_trace_metrics.py} and
  \path{Paper/revision_2026-06-19/source_trace_metrics.json};
  \item short-chain TensorFlow saved-model rerun:
  \path{Paper/revision_2026-06-19/scripts/recompute_short_chain_nn_metrics.py}
  and \path{Paper/revision_2026-06-19/short_chain_nn_rerun_metrics.json};
  \item eigen-relaxation diagnostic script and output:
  \path{Paper/revision_2026-06-19/scripts/analyze_eigen_fit_windows.py} and
  \path{Paper/revision_2026-06-19/eigen_fit_sensitivity.json};
  \item core numerical claim-audit script:
  \path{Paper/revision_2026-06-19/scripts/audit_manuscript_claims.py};
  \item core numerical claim-audit outputs:
  \path{Paper/revision_2026-06-19/manuscript_claim_audit.json} and
  \path{Paper/revision_2026-06-19/manuscript_claim_audit.md};
  \item data/figure availability path audit:
  \path{Paper/revision_2026-06-19/scripts/audit_availability_paths.py},
  \path{Paper/revision_2026-06-19/availability_path_audit.json}, and
  \path{Paper/revision_2026-06-19/availability_path_audit.md};
  \item BibTeX reference database and integrity-audit snapshot:
  \path{Paper/revision_2026-06-19/references.bib} and
  \path{Paper/revision_2026-06-19/integrity_audit_2026-06-19.md}.
\end{itemize}
\endgroup
A public repository snapshot is maintained at
\href{https://github.com/jayleenjiang/nls-energy-cascade}{the project
repository on GitHub}.

\section*{Declarations}

\paragraph{Ethics approval.}
This study is a numerical and theoretical investigation of stochastic
dynamical systems and does not involve human participants, animal subjects, or
personal data.

\paragraph{Author contributions.}
Jayleen Jiang performed the numerical experiments, assembled the computational
artifacts, and drafted the manuscript. Yao Li supervised the project and
contributed to the model formulation, theoretical framing, and interpretation.
Both authors should review and approve the final submitted version.

\paragraph{Competing interests.}
No competing interests are declared in the materials supplied for this draft.

\paragraph{Funding.}
Funding information was not supplied in the current manuscript materials and
should be completed by the authors before submission if required by the target
journal.

\paragraph{AI-assisted preparation.}
AI assistance was used during revision for code audit, consistency checking,
LaTeX editing, and reproducibility organization under author direction. The
authors remain responsible for verifying the scientific claims, numerical
results, and final wording before submission.


\bibliography{references}


\end{document}
